% Generated draft: requires mathematical and source review.
\documentclass[11pt]{article}
\usepackage[margin=25mm]{geometry}
\usepackage{fontspec}
\setmainfont{Latin Modern Roman}
\usepackage{amsmath,amssymb,mathtools,unicode-math}
\setmathfont{Latin Modern Math}
\usepackage{enumitem,needspace,xurl,hyperref}
\setlength{\parindent}{0pt}
\setlength{\parskip}{5pt}
\setlength{\emergencystretch}{3em}
\newcommand{\N}{\mathbb{N}}
\newcommand{\Z}{\mathbb{Z}}
\newcommand{\Q}{\mathbb{Q}}
\newcommand{\R}{\mathbb{R}}
\newcommand{\C}{\mathbb{C}}
\newcommand{\F}{\mathbb{F}}
\newcommand{\A}{\mathbb{A}}
\newcommand{\PP}{\mathbb{P}}
\newcommand{\E}{\mathbb{E}}
\newcommand{\eps}{\varepsilon}
\newcommand{\dd}{\,\mathrm{d}}
\begin{document}
\section*{UnsolvedMath ID: 800885}
\textbf{Source problem code:} catalogue entry
\subsection{Definitions and mathematical statement}
Let \(n\geq 1\) and let \(\Omega\subset\C^n\) be a bounded pseudoconvex domain. Let \(\mu\) be a positive Borel measure on \(\Omega\) with \(\mu(\Omega)=1\). A function is plurisubharmonic (psh) if it is upper semicontinuous and its restriction to each complex line is subharmonic wherever defined. Use the convention \(dd^c=(i/\pi)\partial\bar\partial\); for locally bounded psh functions, \((dd^c u)^n\) denotes the complex Monge--Amp\`ere measure in the Bedford--Taylor sense.

For a precise global interpretation of H\"older continuity, require a continuous extension to \(\overline\Omega\) and constants \(C<\infty\) and \(0<\alpha\leq 1\) such that
\[
 |u(z)-u(w)|\leq C\lVert z-w\rVert^\alpha
 \qquad(z,w\in\overline\Omega).
\]
The boundary condition means that this extension vanishes on \(\partial\Omega\). The supplied source says H\"older continuous in \(\Omega\), without explicitly specifying whether the estimate is global or merely local; the global interpretation used here therefore requires confirmation.

\textbf{Question 16 (attributed to Zeriahi).} Suppose there exists a H\"older continuous psh function \(u\) with zero boundary values such that
\[
 \mu\leq(dd^c u)^n\quad\text{in }\Omega,
\]
where domination means \(\mu(E)\leq(dd^c u)^n(E)\) for every Borel set \(E\subset\Omega\). Must there exist a H\"older continuous psh function \(\varphi\), also extending continuously with zero values on \(\partial\Omega\), such that
\[
 (dd^c\varphi)^n=\mu\quad\text{in }\Omega?
\]
The H\"older exponent and constant for \(\varphi\) are existential and need not equal those for \(u\); the source prescribes no relation between them. No boundary smoothness or stronger geometric assumption on \(\Omega\) is supplied.
\subsection{Short English statement}
Given a probability measure on a bounded pseudoconvex domain, does a H\"older continuous plurisubharmonic subsolution with zero boundary values guarantee a H\"older continuous plurisubharmonic solution of the complex Monge--Amp\`ere equation with the same boundary values? The precise scope of H\"older continuity requires confirmation. The supplied registry labels this question open; that status has not been independently verified.
\subsection{Sources}
\begin{itemize}
\item Local UnsolvedMath registry, permanent ID \(800885\), source problem code \texttt{catalogue entry}, title: \emph{Open Problems in Pluripotential Theory --- Question 16}. The supplied record attributes the collection to Dinew, Guedj, and Zeriahi (2015), and the individual question to Zeriahi.
\item Dinew, Guedj, and Zeriahi, \emph{Open Problems in Pluripotential Theory} (2015), Question 16. Supplied source URL: \url{https://arxiv.org/abs/1511.00705}. The provenance reports access on 2026-10-16 and extraction by \texttt{arxiv\_latex}; these are supplied provenance claims, not independently verified access or source checking.
\end{itemize}
\end{document}
